% ==============================================================================
% IEEE Transactions / Conference Paper Template (skeleton)
% Based on IEEEtran v1.8b  —  CTAN: https://ctan.org/pkg/ieeetran
% Class IEEEtran is contained in standard TeX Live / MiKTeX, so this file
% compiles out-of-the-box with:  pdflatex -> bibtex -> pdflatex -> pdflatex
%
% Document-class modes (pick one):
%   journal      IEEE Transactions / Journals (default; 10pt, twocolumn)
%   conference   IEEE conference proceedings (10pt, twocolumn, looser margins)
%   lettersize   IEEE Letters
%   technotes    Technical notes / letters
% Class options:
%   10pt/11pt/12pt   letterpaper/a4paper   twoside/oneside
%   final/draft      compsoc (Computer Society style)   amsart (math hooks)
%
% Typography (IEEEtran defaults):
%   Body font  : Times Roman (10pt journal) via OT1/T1 encoding
%   Math font  : Computer Modern math symbols (or mathptmx if enabled)
%   Line spacing: single (1.0); captions 8pt
% Citation style: numeric [1], [2], [3]-[5]  (package: cite)
% Bibliography  : \bibliographystyle{IEEEtran}  (IEEEtran.bst)
% Figures       : single column = \columnwidth (3.5in / 88mm)
%                 double column = \textwidth  or figure*  (7.16in / 180mm)
% Caption       : below for figures, above for tables; 8pt font
% ==============================================================================
\documentclass[10pt,journal,final]{IEEEtran}
% For a conference paper instead:
% \documentclass[conference]{IEEEtran}
% For Computer Society journals:
% \documentclass[compsoc]{IEEEtran}

% ---- Core packages ----
\usepackage[T1]{fontenc}
\usepackage[utf8]{inputenc}
\usepackage{cite}                 % IEEE numeric citation sorting/compression
\usepackage{amsmath,amssymb,amsfonts}
\usepackage{algorithm}            % algorithm floating environment
\usepackage{algorithmic}          % IEEE-style algorithm pseudo-code
\usepackage{graphicx}
\usepackage{textcomp}
\usepackage{xcolor}
\usepackage{booktabs}             % professional table rules
\usepackage{url}                  % \url{} inside bib entries
\usepackage{array}
% Optional (load hyperref LAST if used):
% \usepackage{hyperref}
% \hypersetup{colorlinks=true,linkcolor=red,citecolor=blue,urlcolor=magenta}

% ---- IEEEtran helper commands (provided by the class) ----
% \IEEEPARstart{D}{rop}  drop-cap first letter of the introduction
% \IEEEpubid{...}        publication ID footer
% \IEEEtriggeratref{N}   trigger a column break in the reference list

% ==============================================================================
% Title and authors
% ==============================================================================
\title{Title of the IEEE Transactions Submission:\\
A Two-Line Informative Title}

\author{First~Author,~Second~Author,~and~Senior~Author,~\IEEEmembership{Fellow,~IEEE}%
  \thanks{Manuscript received Month DD, YYYY; revised Month DD, YYYY. (Corresponding author: First Author.)}%
  \thanks{The authors are with the Department of X, University of Y, City, Country (e-mail: first.author@univ.edu).}}

\begin{document}

\maketitle

% ==============================================================================
% Abstract & Index Terms  (IEEE style: <= 200 words, single paragraph, no cites)
% ==============================================================================
\begin{abstract}
% TODO: 150-200 words, single paragraph. State the problem (1-2 sentences),
% the approach (2-3 sentences), the key quantitative results (2-3 sentences),
% and the significance (1 sentence). No citations, no undefined abbreviations.
\end{abstract}

\begin{IEEEkeywords}
% TODO: 3-6 keywords separated by commas (IEEE indexing terms).
keyword one, keyword two, keyword three.
\end{IEEEkeywords}

% ==============================================================================
% I. Introduction
% ==============================================================================
\section{Introduction}
\IEEEPARstart{D}{rop}~the first letter here. % TODO: replace with your real opening word.
% TODO: motivate the problem; summarise prior limitations; state contributions.
% Recommended structure: (i) context, (ii) gap, (iii) approach, (iv) bulleted
% contribution list, (v) paper roadmap. Reference prior work with \cite{ref_example}.

% ==============================================================================
% II. Related Work
% ==============================================================================
\section{Related Work}
% TODO: cluster prior work by approach; compare on methodology, data, performance;
% explicitly position your contribution against each cluster in 1-2 sentences.

% ==============================================================================
% III. Methods / System Model
% ==============================================================================
\section{Methods}
% TODO: problem formulation, notation table, model assumptions, the proposed
% method with algorithm pseudo-code, and complexity analysis.

\begin{algorithm}[t]
\caption{Outline of the proposed method}
\label{alg:method}
\begin{algorithmic}[1]
\State \textbf{Input:} dataset $\mathcal{D}$, hyper-parameters $\Theta$
\State initialise model parameters $\theta$
\While{not converged}
    \State update $\theta \leftarrow \theta - \eta \,\nabla \mathcal{L}(\theta;\mathcal{D})$
\EndWhile
\State \textbf{return} trained model
\end{algorithmic}
\end{algorithm}

% ==============================================================================
% IV. Experiments
% ==============================================================================
\section{Experiments}
% TODO: experimental setup (datasets, baselines, metrics, hardware),
% main results table, ablation study, qualitative examples, statistical tests.

\begin{figure}[t]
\centering
% \includegraphics[width=\columnwidth]{figures/result.pdf}
\fbox{\rule{0pt}{2.5cm}\rule{\columnwidth}{0pt}}
\caption{Single-column figure. Caption goes BELOW the figure. Use
\texttt{\textbackslash columnwidth} for single column; for a double-column
figure use \texttt{figure*} and \texttt{\textbackslash textwidth}.}
\label{fig:result}
\end{figure}

\begin{table}[t]
\caption{Table caption goes ABOVE the table (IEEE convention).}
\label{tab:results}
\centering
\begin{tabular}{lcc}
\toprule
Method   & Metric A       & Metric B       \\
\midrule
Baseline & 0.812          & 0.754          \\
Ours     & \textbf{0.891} & \textbf{0.832} \\
\bottomrule
\end{tabular}
\end{table}

% ==============================================================================
% V. Conclusion
% ==============================================================================
\section{Conclusion}
% TODO: one paragraph recapping the contribution and the headline result,
% followed by one sentence of future work. No new citations or claims here.

% ==============================================================================
% References  —  IEEEtran numeric style [1], [2], [3]-[5]
% Compile with:  pdflatex main  ->  bibtex main  ->  pdflatex main  ->  pdflatex main
% ==============================================================================
\bibliographystyle{IEEEtran}
% Drop your .bib file as references.bib in the same folder:
\bibliography{references}
% --- Inline fallback (uncomment if you do not use BibTeX) ---------------------
% \begin{thebibliography}{99}
% \bibitem{ref_example} A.~Author and B.~Author, ``Title of the paper,''
%   \emph{IEEE Trans. Signal Process.}, vol.~XX, no.~Y, pp.~1--10, Mon. 20YY.
% \end{thebibliography}

\end{document}
